% Backup de la sección 4.2 (Consecuencias para el SAT y subclases tratables)
% del capítulo 4 (capitulo_4_modelacion_sat.tex), tomado ANTES de darle tratamiento.
% Fecha: 2026-06-17. Solo backup; no editar aquí.

\section{Consecuencias para el SAT y subclases tratables}
\label{sec:lectura-sat}

Con la enumeración del capítulo~\ref{cap:vacio-interseccion-rcg-regular}
se decide esa satisfacibilidad: se
recorren las cadenas aceptadas por el autómata booleano y se aplica a
cada una el reconocimiento de la RCG de consistencia. Las dos partes
de la instancia tienen costo polinomial. El autómata booleano tiene
$O(n^2)$ estados~\cite{fernandezarias2007}, con $n$ la cantidad de
ocurrencias de variables booleanas en la fórmula. La gramática
$G_{\mathrm{cons}}$ tiene aridad uno y tamaño $O(nm)$, contado en
símbolos, con $m$ la cantidad de variables booleanas distintas: dos
cláusulas de longitud $O(n)$ por variable booleana, más las $O(n)$
cláusulas de los predicados $L_j$. El reconocimiento con
$G_{\mathrm{cons}}$ es polinomial: cada cláusula admite a lo sumo una sustitución de rango
para cada cadena, porque los rangos de las variables $W_k$ quedan
determinados por las longitudes $g_k$, y esa sustitución se comprueba
en tiempo lineal; el reconocimiento de una cadena de longitud $n$
cuesta entonces $R_{G_{\mathrm{cons}}}(n) = O(nm)$.

Toda cadena aceptada por el autómata booleano tiene longitud $n$, de
modo que cada llamada a $\proc{RCG-Recognize}$ cuesta $O(nm)$. Con
$|Q| = O(n^2)$, la cota de la
subsección~\ref{sec:complejidad-algoritmo} queda en
\[
  O\bigl( n^2 + p(A)\,(n^2 + nm) \bigr) = O\bigl( n^2 + p(A)\,n^2 \bigr),
\]
pues $m \le n$. El factor dominante es $p(A)$, la cantidad de cadenas
aceptadas por el autómata booleano, que para una fórmula con muchas
interpretaciones verdaderas bajo el supuesto de independencia es del
orden de $2^n$. La dificultad no está entonces en el tamaño de la
instancia ni en el reconocimiento, sino en el vacío: decidir si
$L(G_{\mathrm{cons}}) \cap L(A)$ es vacío es el SAT, y para una RCG
con un autómata finito acíclico ese problema es
NP-completo~\cite{bertsch2001rcg}.

En SATCF y SATSRC ese costo no aparece, y la razón no es una enumeración
más hábil. En esos enfoques el vacío se decide desde el formalismo: el
de una gramática libre del contexto y el de una sRCG se deciden en
tiempo polinomial, sin recorrer las
interpretaciones~\cite{fernandezarias2007,aguileralopez2016}. Para las
RCG no hay ese procedimiento directo: con la enumeración el problema se
decide, pero el costo es el número de interpretaciones, no una propiedad
del formalismo.

La no-linealidad tiene dos efectos. Por un lado, con la igualdad
entre las subcadenas en que aparece una misma variable booleana se amplía la
clase de fórmulas cuya consistencia se expresa, más allá del orden de
concatenación de rango simple de las sRCG. Por otro, por esa misma
igualdad el vacío de la intersección deja de ser polinomial.

Quedan entonces dos vías hacia una extensión tratable de SATSRC. En
la primera se restringen las fórmulas: si el autómata booleano tiene
una cantidad de caminos aceptantes acotada por un polinomio en $n$,
la cota anterior es polinomial y el algoritmo de enumeración
decide la satisfacibilidad en tiempo polinomial, por lo que ya esa
propuesta queda cerrada.

En la segunda vía se restringe la gramática de consistencia a una
subclase no lineal de las RCG, que debe conservar parte de la
expresividad que la no-linealidad añade sobre las sRCG y admitir un
procedimiento polinomial para el vacío de la intersección que no
recorra las interpretaciones. Con esta vía también se resuelve solo
una clase de fórmulas, aquellas cuyas cadenas consistentes se
reconocen por una gramática de la subclase: igual que en SATCF y
SATSRC, la clase de fórmulas queda delimitada por el formalismo de la
gramática de consistencia. Para hallar esa subclase hace falta una
representación simbólica de la intersección, de tamaño polinomial
respecto a la fórmula, sobre la que estudiar con qué restricciones el
vacío se vuelve tratable.

En el capítulo siguiente se construye esa representación: una gramática
producto entre una RCG y un lenguaje regular. Con ella el problema
general no se vuelve polinomial, pero se obtiene el objeto sobre el que
estudiar las subclases no lineales con vacío polinomial.
